\section{优化基础}

既然我们知道了如何衡量一组权重 $W$ 的好坏（通过损失函数），下一个问题是：如何找到最优的 $W$？

\subsection{损失景观与优化直觉}

\begin{figure}[H]
\centering
\includegraphics[width=0.95\textwidth]{slides-images/slide-020.jpg}
\caption{损失景观：纵轴为损失值，水平面为参数空间。目标是找到景观中的最低点\protect\footnotemark}
\end{figure}
\footnotetext{Slides 第24页。}

可以将优化想象为一个人在损失景观中行走，试图找到最低的山谷。关键的类比是：这个人是\textbf{蒙着眼的}——他看不到远处的地形，只能感受脚下地面的坡度，然后决定向哪个方向走。

\subsection{策略一：随机搜索}

最简单的方法是随机尝试大量不同的 $W$，选择损失最小的那个：

\begin{lstlisting}[caption={随机搜索：尝试 1000 个随机权重}]
best_loss = float('inf')
for i in range(1000):
    W = np.random.randn(10, 3072) * 0.0001
    loss = L(X_train, Y_train, W)
    if loss < best_loss:
        best_loss = loss
        best_W = W
\end{lstlisting}

这种方法在 CIFAR-10 上只能达到约 15.5\% 的准确率（随机猜测为 10\%），而现代深度学习方法可以达到 99.7\%。显然不够好，但胜于随机猜测。

\subsection{策略二：跟随坡度——梯度}

更好的策略是``感受地面的坡度，朝下坡方向走''。数学上，这就是计算\textbf{梯度}。

\textbf{一维导数}的定义（极限形式）：

$$\frac{df}{dx} = \lim_{h \to 0} \frac{f(x + h) - f(x)}{h}$$

\begin{figure}[H]
\centering
\includegraphics[width=0.95\textwidth]{slides-images/slide-024.jpg}
\caption{一维导数与多维梯度：导数给出函数在当前点的变化率\protect\footnotemark}
\end{figure}
\footnotetext{Slides 第28页。}

在多维情况下，使用\textbf{梯度}——对每个维度分别计算偏导数，得到一个向量：

$$\nabla_W L = \left[\frac{\partial L}{\partial W_1}, \frac{\partial L}{\partial W_2}, \ldots\right]$$

梯度指向函数增长最快的方向，因此\textbf{负梯度方向}是下降最快的方向。

\subsection{数值梯度 vs 解析梯度}

\begin{figure}[H]
\centering
\includegraphics[width=0.95\textwidth]{slides-images/slide-025.jpg}
\caption{数值梯度的计算：对每个参数加一个很小的 $h$，计算损失变化\protect\footnotemark}
\end{figure}
\footnotetext{Slides 第29页。}

\textbf{数值梯度}：用有限差分近似导数（$h$ 取一个很小的值如 $10^{-5}$）
\begin{itemize}
    \item 优点：实现简单
    \item 缺点：速度慢（需要遍历每个参数）、近似值（不精确）
\end{itemize}

\textbf{解析梯度}：通过微积分推导损失函数对 $W$ 的导数公式
\begin{itemize}
    \item 优点：精确、快速
    \item 缺点：推导可能出错
\end{itemize}

\begin{importantbox}{梯度检查（Gradient Check）}
实践中的标准做法是：用解析梯度进行训练（快速且精确），但用数值梯度进行\textbf{验证}——确保两者结果一致。如果解析梯度的实现有 bug，梯度检查可以立即发现。这是作业中常见的调试手段。
\end{importantbox}

\subsection{本章小结}

优化的目标是在损失景观中找到最低点。梯度给出了损失函数在当前点的下降方向。实践中使用解析梯度计算，并通过数值梯度进行验证。接下来将介绍具体的优化算法。

%% ========== 梯度下降及其变体 ==========

